\section{Reading Note: ABR09}

\subsection{Dependence Structure}

\subsection{Proof of Corollary 10 from Lemma 9}
We explicitly write out the proof sketched in the paper. The main ideas used is the ``amplification argument'', which is considering independent sub BRWs from vertices of the original BRW. 

\begin{tcolorbox}[title=,colback=gray!10,colframe=black,breakable]
    Lemma 9. For all values of $m$ for which $m^*-m=o(\sqrt{n})$,

\[
\mathbf{E}\left\{\left|G_{n, m} \|\right| v_n \in G_{n, m}\right\}=\mathbf{E}\left|G_{n, m}\right|+O(1)
\]


Assume for the moment that Lemma 9 holds, and let

\[
m^{\prime}=m^*-\frac{\log n}{2 t}=\Lambda^{\prime}(t) n-\frac{3 \log n}{2 t}
\]


By Lemma 7 and linearity of expectation, we see that $m^{\prime}$ is within $O(1)$ of the smallest value $m$ for which $\mathbf{E}\left|G_{n, m}\right| \geq 1$. Applying (25) with $\mathcal{R}=G_{n, m^{\prime}}$, it follows immediately that $\mathbf{P}\left\{\left|G_{n, m^{\prime}}\right|>0\right\}=\Omega(1)$. Since $G_{n, m} \subseteq N_{n, m}$ for all $m$, it follows that $\mathbf{P}\left\{\left|N_{n, m^{\prime}}\right|>0\right\}=\Omega(1)$, so $\mathbf{P}\left\{M_n \leq m^{\prime}\right\}=\Omega(1)$. By applying an amplification argument exactly as we did when proving Theorem 4, we immediately obtain exponential tail bounds for the upper tail of $M_n$ :

Corollary 10. There exist constants $C_1>0, \delta_1>0$ such that for all $x>0$,

\[
\mathbf{P}\left\{M_n \geq m^{\prime}+x \mid \delta\right\} \leq C_1 e^{-\delta_1 x}
\]

\end{tcolorbox}

\begin{proof}[Proof of Corollary~10]
    \textbf{Step 1.} We have previously seen, by calculation using Theorem~1, that $m'$ is within $O(1)$ of $\Gamma'(t) n - 3 \log n / 2t$, the smallest value of $m$ such that $P(S_v \le m) \ge n / (EB)^n$. Recall the definition of $G_{n, m}$:
    \[
G_{n, m}=\left\{v \in N_n: m-1 \leq S_v \leq m, S_v \text { is strictly leading }\right\},
\]
    and note that $E |G_{n, m}| = |N_{n}| P(S_v \in G_{n, m})$. By Lemma~7, we have $P(S_v \in G_{n, m}) = \Theta(P(S_v \le m) / n).$ Thus, $E |G_{n, m}| = \Theta(1)$.

    \textbf{Step 2.} We apply the Chung-Erdős inequality:
    \begin{tcolorbox}[title=Chung-Erdős Inequality,colback=gray!10,colframe=black,breakable]
        for any integer $i \geq 1$ and any random set $\mathcal{R} \subseteq N_n^{\prime}$ (recall that $N_n^{\prime}$ is the set of nodes at depth $n$ in $T_n^{\prime}$ ),
\[
\mathbf{P}\{|\mathcal{R}|>0\} \geq \frac{\mathbf{E}|\mathcal{R}|}{1+\sup _{v \in N_n^{\prime}} \mathbf{E}\{| \mathcal{R} | \mid  v \in \mathcal{R}\}}
\]

    Below is a format statement in its original form:

    Let $A_1, \ldots, A_n$ be events. Assume that $\operatorname{Pr}\left[A_i\right]>0$ for some $i$. Then

    \[
    \operatorname{Pr}\left[A_1 \vee \cdots \vee A_n\right] \geq \frac{\left(\sum_{i=1}^n \operatorname{Pr}\left[A_i\right]\right)^2}{\sum_{i=1}^n \sum_{j=1}^n \operatorname{Pr}\left[A_i \wedge A_j\right]}
    \]

    The proof is by Paley-Zygmund inequality.
\end{tcolorbox}

Take $\mathcal{R} = G_{n, m'}$. Since $|G_{n, m'}| > 1$ for our choice of $m'$, we have $\mathbf{P}\left\{|G_{n, m'}| > 0\right\} = \Omega(1)$. Since $G_{n, m'} \subseteq N_{n, m'}$ for all $m$, it follows that $\mathbf{P}\left\{|N_{n, m'}| > 0\right\} = \Omega(1)$. Thus, $\mathbf{P}\left\{M_n \leq m'\right\} = \Omega(1)$.

\textbf{Step 3.} We now apply an amplification argument. Intuitively, we first identify some vertices to the left of $x$ at some intermediate time. For the event $\{M_n \geq m' + x, S\}$ to happen, their subtrees must all travel at least $m'$ steps to the right, which is unlikely due to Step 2. 

Since the process is supercritical, we apply the following inequality due to (McDiarmid, 1995, Lemma 1) to control the number of vertices from below:
\[
\mathbf{P}\left\{0<\left|N_i\right| \leq \gamma_0^i\right\} \leq c_0 e^{-\delta_0 i}\]
\textbf{Intuition.} $|N_i|$ is expected to grow like $B^i$, thus we expect $E |N_i| \sim \exp(i \log EB)$. Choosing a smaller exponent $\gamma_0$ will result in an exponentially decaying tail bound.

Take $\ell(x) = \lceil \log _{\gamma_0} x \rceil$. Then, for $n \ge \ell(x)$, we have
\begin{align*}
    \begin{aligned}
    \mathbf{P}\left\{M_n>x + m', \mathscr{S}\right\} 
    & \leq \mathbf{P}\left\{0<\left|N_{\lfloor\ell(x)\rfloor}\right| \leq x\right\}+\mathbf{P}\left\{M_n>x + m' | | N_{\ell(x)} \mid>x\right\} \\
    & \leq c_0 e^{-\delta_0 x}+\sup _{x<k \leq d^{\ell(x)}} \mathbf{P}\left\{M_n>x + m'| | N_{\ell(x)} \mid=k\right\} .
    \end{aligned}
\end{align*}
Let $F_{x, d}$ be the set of nodes with depth $d$ and position at most $x$.
For fixed $k$ in the above range, we have

\[
\begin{aligned}
& \mathbf{P}\left\{M_n>x + m'| | N_{\ell(x)} \mid=k\right\} \\
& =\sum_{i=0}^k\left(\mathbf { P } \left\{M_n>x + m' | | N_{\ell(x)}\left|=k,\left|F_{x, \ell(x)}\right|=i\right\}\right.\right. \\
& \left.\cdot \mathbf{P}\left\{\left|F_{x, \ell(x)}\right|=i \| N_{\ell(x)} \mid=k\right\}\right) \\
& \leq \sum_{i=0}^k\left(1-p_0\right)^i \mathbf{P}\left\{\left|F_{x, \ell(x)}\right|=i \| N_{\ell(x)} \mid=k\right\} \\
& \leq\left(1-p_0\right)^{k / 2}+\mathbf{P}\left\{\left|F_{x, \ell(x)}\right| \leq k / 2| | N_{\ell(x)} \mid=k\right\} ,
\end{aligned}
\]
where $p_0 = \mathbf{P}\left\{M_n \geq m' \right) = \Omega(1)$.

The rest follows word-by-word. 

\textbf{Intuition}. Since $\ell(x)$ grows logarithmically with $x$, for a vertice in $N_{\ell(x)}$ to be outside of $F_{x, \ell(x)}$, it has to travel exponentially many steps to the right, which is unlikely. Thus the bound of the right hand side.
\end{proof}
To summarize again: We first run the walk for $\ell(x)$ steps to obtain around $k$ vertices. It is very unlikely for most of them to go beyond $x$, since it requires exponential speedup. On the other hand, for $M_n \ge x + m'$, the subtree of each of them must travel at least $m'$ steps to the right, which has a probability strictly less than 1. Thus, the probability of the event is exponentially small in $x$.


\subsection{Proof of Lemma 7}

We first give a hand-wavy yet intuitive proof of Lemma 6, and then demonstrate how it implies Lemma 7.

\begin{tcolorbox}[title=,colback=gray!10,colframe=black,breakable]
    LEMMA 6. Given a random walk $S$ with steps distributed as $X$, and real numbers $a, c$ with $c>\mathbf{E} X \geq a$ and for which $\mathbf{P}\{\mathbf{E} X<X \leq c\}>0$,

\[
\mathbf{P}\left\{S_n \leq a n, S_n \text { is strictly leading }\right\} \leq \frac{\mathbf{P}\left\{S_n \leq a n\right\}}{n}
\]

and

\[
\mathbf{P}\left\{S_n \leq a n, S_n \text { is strictly leading }\right\} \geq \mathbf{P}\{\mathbf{E} X<X \leq c\} \cdot \frac{\mathbf{P}\left\{S_{n-1} \leq a n-c\right\}}{n-1}
\]

\end{tcolorbox}

For the upper bound: (1) By drawing a picture, one is convinced that at most one path in all cyclic permutations can be strictly leading. Draw it like a periodic function, and slide the window to represet a permuted walk. (2) The event $S_n \le an$ is fixed by all cyclic permutations.

For the lower bound: We showcase a strictly leading path among the permutations. First, note that for $S_n \le a n$ and $S_n$ strictly leading, it suffices to have the first step $X_1$ greater than $EX_1$ and the rest of the steps leading, with appropriate condition on the last step. Then, one sees that \textit{any} path has at least one leading path among its cyclic permutations.


\subsection{Proof of Lemma 9} 

For this one we need two controls on probability of $S_n$ staying above a level $a$ and ending at some site $m$. In the proofs we use Lemma 7, the cyclic permutation, and Theorem 1, the large deviation estimate.

